% ============================================================
%  BID-IA — Agente de Revisión de Literatura
%  Presentación para revisión interna — Fedesarrollo & BID
%  Mayo 2026
% ============================================================
\documentclass[aspectratio=169,10pt]{beamer}

\usepackage[utf8]{inputenc}
\usepackage[T1]{fontenc}
\usepackage[spanish]{babel}
\usepackage{lmodern}
\usepackage{tikz}
\usepackage{pgfplots}
\pgfplotsset{compat=1.18}
\usepackage{booktabs}
\usepackage{colortbl}
\usepackage{xcolor}
\usepackage{fontawesome5}
\usepackage{tcolorbox}
\usepackage{array}
\tcbuselibrary{most}
\usetikzlibrary{arrows.meta,positioning,shapes.geometric,fit,backgrounds,calc}

% ── Colores ───────────────────────────────────────────────────
\definecolor{bidblue}{HTML}{00296B}
\definecolor{bidmid}{HTML}{003F88}
\definecolor{bidlight}{HTML}{00509D}
\definecolor{bidaccent}{HTML}{4AABE8}
\definecolor{bidgray}{HTML}{CBD5E1}
\definecolor{colgreen}{HTML}{16A34A}
\definecolor{colamber}{HTML}{D97706}
\definecolor{colred}{HTML}{DC2626}
\definecolor{darkbg}{HTML}{0A1628}
\definecolor{cardbg}{HTML}{0F1F3D}

% ── Tema ──────────────────────────────────────────────────────
\usetheme{default}
\usecolortheme{default}
\setbeamercolor{normal text}{fg=white,bg=bidblue}
\setbeamercolor{structure}{fg=bidaccent}
\setbeamercolor{frametitle}{fg=white,bg=bidmid}
\setbeamercolor{title}{fg=white}
\setbeamercolor{block title}{fg=white,bg=bidlight}
\setbeamercolor{block body}{fg=white,bg=cardbg}
\setbeamerfont{frametitle}{size=\large,series=\bfseries}
\setbeamerfont{title}{size=\LARGE,series=\bfseries}
\setbeamertemplate{navigation symbols}{}
\setbeamertemplate{footline}{%
  \leavevmode
  \hbox{%
    \begin{beamercolorbox}[wd=.5\paperwidth,ht=2.25ex,dp=1ex,leftskip=4mm]{author in head/foot}%
      \color{bidaccent!60}\tiny BID-IA \textperiodcentered\ Fedesarrollo%
    \end{beamercolorbox}%
    \begin{beamercolorbox}[wd=.5\paperwidth,ht=2.25ex,dp=1ex,rightskip=4mm,right]{date in head/foot}%
      \color{bidaccent!60}\tiny\insertframenumber\,/\,\inserttotalframenumber%
    \end{beamercolorbox}}%
  \vskip0pt}
\setbeamertemplate{frametitle}{%
  \nointerlineskip
  \begin{beamercolorbox}[wd=\paperwidth,ht=2.8ex,dp=1.2ex,leftskip=8mm]{frametitle}%
    \usebeamerfont{frametitle}\insertframetitle
    \ifx\insertframesubtitle\empty\else
      \hspace{6pt}\textcolor{bidaccent!70}{\normalfont\small\insertframesubtitle}%
    \fi
  \end{beamercolorbox}}

% ── Helpers ───────────────────────────────────────────────────
\newcommand{\hl}[1]{\textcolor{bidaccent}{\textbf{#1}}}

\newcommand{\statuspill}[2][colgreen]{%
  \tikz[baseline=(n.base)]\node[fill=#1!80!black,rounded corners=2pt,
    inner xsep=4pt,inner ysep=1pt,font=\tiny\bfseries,text=white](n){#2};\relax}

\newenvironment{card}[1][cardbg]{%
  \begin{tcolorbox}[colback=#1,colframe=bidlight!40,arc=4pt,
    boxrule=0.4pt,left=5pt,right=5pt,top=4pt,bottom=4pt]}%
  {\end{tcolorbox}}

\newenvironment{qbox}[1]{%
  \begin{tcolorbox}[colback=bidlight!25,colframe=bidaccent,arc=4pt,boxrule=0.6pt,
    title={\small\bfseries #1},fonttitle=\color{white},coltitle=white,
    attach boxed title to top left={yshift=-2mm,xshift=3mm},
    boxed title style={colback=bidaccent,arc=3pt},
    left=5pt,right=5pt,top=6pt,bottom=4pt]}%
  {\end{tcolorbox}}

% ============================================================
\begin{document}
% ============================================================

% ── Portada ──────────────────────────────────────────────────
{
\setbeamercolor{background canvas}{bg=bidblue}
\begin{frame}[plain]
\vspace{10mm}\hspace{7mm}%
\begin{minipage}{0.76\textwidth}
  {\tiny\color{bidaccent}\bfseries FEDESARROLLO \& BANCO INTERAMERICANO DE DESARROLLO}
  \vspace{4mm}

  {\huge\bfseries\color{white}Agente de Revision\\[3pt]de Literatura con IA}
  \vspace{5mm}

  {\large\color{bidaccent}BID-IA · Adopcion de IA en Firmas — Colombia y LATAM}
  \vspace{6mm}

  \textcolor{bidgray!80}{\small
    Arquitectura, metodologia de evaluacion\\
    y resultados de la primera iteracion.}
  \vspace{8mm}

  \textcolor{bidgray!50}{\footnotesize Mayo 2026 \quad·\quad Maria Alejandra Montoya}
\end{minipage}
\end{frame}
}

% ── Índice ────────────────────────────────────────────────────
\begin{frame}{Contenido}
\begin{columns}[T]
\column{0.48\textwidth}
\begin{enumerate}\setlength\itemsep{6pt}
  \item[\textcolor{bidaccent}{1.}] \textbf{El problema} --- ¿por que un agente?
  \item[\textcolor{bidaccent}{2.}] \textbf{Arquitectura} --- como esta construido
  \item[\textcolor{bidaccent}{3.}] \textbf{El loop agentico} --- como decide
  \item[\textcolor{bidaccent}{4.}] \textbf{Fuentes de busqueda}
\end{enumerate}
\column{0.48\textwidth}
\begin{enumerate}\setlength\itemsep{6pt}
  \item[\textcolor{bidaccent}{5.}] \textbf{Sistema de scoring v2}
  \item[\textcolor{bidaccent}{6.}] \textbf{Resultados} --- primera iteracion
  \item[\textcolor{bidaccent}{7.}] \textbf{Proximos pasos}
\end{enumerate}
\end{columns}
\vspace{7mm}
\begin{card}
\small\color{bidgray}
\textbf{\color{bidaccent}Objetivo de esta revision:} validar la arquitectura y criterios antes de correr los 6 pockets restantes. Se busca feedback sobre (1)~pertinencia del scoring, (2)~cobertura de fuentes, y (3)~prioridades para revision manual.
\end{card}
\end{frame}

% ══════════════════════════════════════════════════════════════
% 1. EL PROBLEMA
% ══════════════════════════════════════════════════════════════
{
\setbeamercolor{background canvas}{bg=bidmid}
\begin{frame}[plain]\vspace{20mm}
\begin{center}
{\Large\bfseries\color{bidaccent}1 · El problema}\\[5pt]
{\color{bidgray}\normalsize ¿Por que necesitamos un agente para la revision de literatura?}
\end{center}
\end{frame}
}

\begin{frame}{El problema}{Escala + criterios simultaneos}
\begin{columns}[T]
\column{0.48\textwidth}
\begin{card}
\textbf{\color{bidaccent}La tarea sin agente}\\[3pt]
\scriptsize\color{bidgray}
\begin{itemize}\setlength\itemsep{2pt}
  \item 7 temas (pockets) $\times$ 60 papers = \hl{420 papers}
  \item Leer abstract, calificar metodologia, verificar relevancia
  \item Decidir: Aceptado / Revisar / Rechazado
  \item Extraer sintesis estructurada de los aceptados
  \item Construir red de conexiones entre papers
\end{itemize}
\vspace{2pt}
\textcolor{bidaccent!60}{\footnotesize $\approx$3--4 semanas por investigador}
\end{card}

\column{0.48\textwidth}
\begin{card}
\textbf{\color{bidaccent}Lo que hace el agente}\\[3pt]
\scriptsize\color{bidgray}
\begin{itemize}\setlength\itemsep{2pt}
  \item Busca en \hl{9 fuentes academicas} simultaneamente
  \item Aplica rubrica de 5 dimensiones de forma consistente
  \item Sintetiza findings, metodologia y relevancia para LATAM
  \item Guarda resultados localmente --- segunda corrida cuesta $\approx$\$0
  \item El investigador solo decide los casos \textbf{REVISAR}
\end{itemize}
\vspace{2pt}
\textcolor{bidaccent!60}{\footnotesize Costo total: \textbf{\$5--7 USD} por corrida completa}
\end{card}
\end{columns}

\vspace{3mm}
\begin{card}
\scriptsize\color{bidgray}
\textbf{\color{bidaccent}Contexto:} esta revision informa el diseno del \textbf{primer RCT de adopcion de IA en MiPyMEs de LATAM} (Guatiguara). La literatura debe responder preguntas concretas: que medir, como identificar efectos causales, que heterogeneidades anticipar.
\end{card}
\end{frame}

% ══════════════════════════════════════════════════════════════
% 2. ARQUITECTURA
% ══════════════════════════════════════════════════════════════
{
\setbeamercolor{background canvas}{bg=bidmid}
\begin{frame}[plain]\vspace{20mm}
\begin{center}
{\Large\bfseries\color{bidaccent}2 · Arquitectura}\\[5pt]
{\color{bidgray}\normalsize ¿Como esta construido el agente?}
\end{center}
\end{frame}
}

\begin{frame}{Arquitectura del sistema}{Tres capas con responsabilidades claras}

\begin{columns}[T]
\column{0.52\textwidth}
\begin{center}
\begin{tikzpicture}[
  every node/.style={font=\scriptsize},
  layer/.style={rectangle,rounded corners=4pt,text=white,align=center,
    minimum width=58mm,minimum height=10mm,inner xsep=5pt,inner ysep=4pt},
  sub/.style={rectangle,rounded corners=2pt,text=white,align=center,
    minimum width=27mm,minimum height=8mm,font=\tiny,inner xsep=3pt,inner ysep=2pt},
  arr/.style={-{Stealth[length=4pt]},bidaccent,line width=0.9pt},
  lbl/.style={font=\tiny,text=bidgray},
]
% Capa 1
\node[layer,fill=bidlight] (cli) at (0,0)
  {\textbf{Capa 1 · Entrada}\\[-1pt]{\tiny\color{bidgray}\texttt{agent.py --pocket X --mode full}}};

% Capa 2
\node[layer,fill=bidblue,minimum height=18mm] (brain) at (0,-1.6)
  {};
\node[font=\tiny\bfseries,text=white] at (0,-1.1) {Capa 2 · Inteligencia};
\node[sub,fill=bidmid] (son) at (-1.5,-1.8) {\textbf{Sonnet 4.6}\\Orquestador};
\node[sub,fill=cardbg,draw=bidaccent!50,line width=0.4pt] (haiku) at (1.5,-1.8) {\textbf{Haiku 4.5}\\Review + sintesis};

% Capa 3
\node[layer,fill=darkbg,draw=bidaccent!30,line width=0.4pt,minimum height=18mm] (data) at (0,-3.3)
  {};
\node[font=\tiny\bfseries,text=bidgray!80] at (0,-2.85) {Capa 3 · Datos y APIs};
\node[sub,fill=colgreen!20!darkbg] (apis) at (-1.5,-3.5) {9 APIs academicas\\SS · NBER · ArXiv · ...};
\node[sub,fill=colamber!15!darkbg] (out) at (1.5,-3.5) {Corpus resultante\\pocket\_reviewed.json};

% Arrows
\draw[arr] (cli)--(brain);
\draw[arr] (son.south)--++(0,-0.45);
\draw[arr] (haiku.south)--++(0,-0.45);
\draw[arr] (apis.north)--++(0,0.45);
\draw[arr] (out.north)--++(0,0.45);
\draw[arr,dashed,bend left=25] (son.north east) to
  node[right,lbl]{historial} (son.north east |- cli.south);
\end{tikzpicture}
\end{center}

\column{0.44\textwidth}
\vspace{2mm}
\begin{card}
\scriptsize\color{bidgray}
\textbf{\color{bidaccent}Capa 1 — Entrada}\\[3pt]
El investigador define el pocket y el modo (\texttt{search}, \texttt{review}, \texttt{full}). Desde ahi, el agente opera de forma autonoma.
\end{card}\vspace{4pt}

\begin{card}
\scriptsize\color{bidgray}
\textbf{\color{bidaccent}Capa 2 — Inteligencia (dos modelos)}\\[3pt]
\textbf{Sonnet} ve el historial completo: planifica, decide que buscar, interpreta resultados y escribe el reporte de gaps.\\[3pt]
\textbf{Haiku} recibe solo el paper a evaluar: aplica la rubrica y sintetiza. 10x mas barato por tarea puntual.
\end{card}\vspace{4pt}

\begin{card}
\scriptsize\color{bidgray}
\textbf{\color{bidaccent}Capa 3 — Datos}\\[3pt]
Las tools Python son deterministas: llaman APIs, enriquecen con Crossref, deduplicaN y guardan. Sonnet nunca toca archivos directamente.
\end{card}
\end{columns}
\end{frame}

\begin{frame}{Arquitectura del sistema}{Que ve y que hace cada componente}
\begin{columns}[T]

\column{0.48\textwidth}
\begin{card}
\scriptsize\textbf{\color{bidaccent}Sonnet 4.6 — Orquestador}\\[3pt]
\tiny\color{bidgray}
\textbf{Recibe en cada turno:}
\begin{itemize}\setlength\itemsep{1pt}
  \item Definicion del pocket (queries, anchor papers, criterios)
  \item Estado actual del corpus (cuantos papers, scores, gaps)
  \item Resultado de la ultima tool ejecutada
  \item Historial de decisiones del loop (max 30 iter.)
\end{itemize}
\vspace{3pt}
\textbf{Decide:} que tool llamar, con que parametros, cuando parar.\\
\textbf{Nunca:} escribe archivos, llama APIs ni evalua papers directamente.
\end{card}

\vspace{3pt}
\begin{card}
\scriptsize\textbf{\color{bidaccent}Haiku 4.5 — Evaluador y Sintetizador}\\[3pt]
\tiny\color{bidgray}
\textbf{Recibe por paper:}
\begin{itemize}\setlength\itemsep{1pt}
  \item Titulo, abstract, metadata enriquecida (DOI, citas, venue)
  \item Definicion del pocket en el \emph{system prompt}
  \item Rubrica de 5 dimensiones (D1--D5)
\end{itemize}
\vspace{3pt}
\textbf{Devuelve:} scores D1--D5, clasificacion (ACEPTADO/REVISAR/RECHAZADO), justificacion por dimension.\\
\textbf{Costo:} $\approx$\$0.0003 USD por paper.
\end{card}

\column{0.48\textwidth}
\begin{card}
\scriptsize\textbf{\color{bidaccent}Tools Python — Las 6 acciones disponibles}\\[3pt]
\tiny\color{bidgray}
\setlength\tabcolsep{3pt}
\renewcommand{\arraystretch}{1.25}
\begin{tabular}{@{}p{23mm}p{34mm}@{}}
\texttt{search\_papers}     & Llama las 9 APIs con las queries del pocket \\
\texttt{enrich\_metadata}   & Crossref: agrega DOI, ISSN, citas formales, ORCID \\
\texttt{deduplicate}        & Merge por DOI + similitud de titulo ($>$0.85) \\
\texttt{evaluate\_coverage} & Score 0--1: papers aceptados vs.\ objetivo \\
\texttt{review\_papers}     & Llama Haiku por cada paper pendiente \\
\texttt{save\_results}      & Merge al JSON del pocket; genera reporte de gaps \\
\end{tabular}
\end{card}

\vspace{3pt}
\begin{card}
\tiny\color{bidgray}
\textbf{\color{bidaccent}Principio de diseno:} Sonnet razona, las tools actuan. Cada decision de Sonnet queda en el historial y cada accion de una tool es reproducible de forma independiente.
\end{card}

\end{columns}
\end{frame}

% ══════════════════════════════════════════════════════════════
% 3. EL LOOP
% ══════════════════════════════════════════════════════════════
{
\setbeamercolor{background canvas}{bg=bidmid}
\begin{frame}[plain]\vspace{20mm}
\begin{center}
{\Large\bfseries\color{bidaccent}3 · El loop agentico}\\[5pt]
{\color{bidgray}\normalsize ¿Como decide el agente en cada iteracion?}
\end{center}
\end{frame}
}

\begin{frame}{El loop agentico}{Secuencia de una corrida --mode full}
\begin{columns}[T]

\column{0.44\textwidth}
\begin{tikzpicture}[
  stepbox/.style={rectangle,rounded corners=3pt,minimum width=46mm,minimum height=8mm,
    text=white,align=center,font=\scriptsize\bfseries,inner sep=3pt},
  arr/.style={-{Stealth[length=3pt]},bidaccent,line width=0.8pt},
  lbl/.style={font=\tiny,text=bidgray},
]
\node[stepbox,fill=bidlight]            (s1) at (0, 0.00) {1 · get\_status + definicion};
\node[stepbox,fill=bidmid]              (s2) at (0,-1.00) {2 · search\_papers x fuentes};
\node[stepbox,fill=bidmid]              (s3) at (0,-2.00) {3 · enrich + deduplicate};
\node[stepbox,fill=cardbg,draw=bidaccent!60] (s4) at (0,-3.00) {4 · evaluate\_coverage};
\node[stepbox,fill=bidblue!80]          (s5) at (0,-4.00) {5 · review with Haiku};
\node[stepbox,fill=bidblue!80]          (s6) at (0,-5.00) {6 · synthesize with Haiku};
\node[stepbox,fill=colgreen!50!bidblue] (s7) at (0,-6.00) {7 · save + reporte};

\draw[arr] (s1)--(s2);
\draw[arr] (s2)--(s3);
\draw[arr] (s3)--(s4);
\draw[arr] (s4)--node[right,lbl]{$\geq$0.6}(s5);
\draw[arr] (s5)--(s6);
\draw[arr] (s6)--(s7);
\draw[arr] (s4.east)--++(3mm,0)--++(0,1.00cm)--++(-3mm,0)
  node[above right,lbl,xshift=-3mm]{$<$0.6};
\end{tikzpicture}

\column{0.52\textwidth}
\begin{card}
\small\textbf{\color{bidaccent}Paso 1 — Estado}\\
\scriptsize\color{bidgray}
Sonnet lee corpus actual y definicion del pocket (queries, anchor papers). Sabe que existe y que falta.
\end{card}\vspace{2pt}
\begin{card}
\small\textbf{\color{bidaccent}Pasos 2--4 — Busqueda}\\
\scriptsize\color{bidgray}
Busca en multiples fuentes, enriquece con Crossref, elimina duplicados. Score de cobertura decide si buscar mas.
\end{card}\vspace{2pt}
\begin{card}
\small\textbf{\color{bidaccent}Pasos 5--6 — IA (Haiku)}\\
\scriptsize\color{bidgray}
Haiku recibe la definicion del pocket en el \emph{system prompt}. Solo sintetiza papers ACEPTADOS.
\end{card}\vspace{2pt}
\begin{card}
\small\textbf{\color{bidaccent}Paso 7 — Guardado}\\
\scriptsize\color{bidgray}
Merge inteligente: agrega papers nuevos sin pisar revisados. Reporte con gaps detectados.
\end{card}
\end{columns}
\end{frame}

% ══════════════════════════════════════════════════════════════
% 4. FUENTES
% ══════════════════════════════════════════════════════════════
{
\setbeamercolor{background canvas}{bg=bidmid}
\begin{frame}[plain]\vspace{20mm}
\begin{center}
{\Large\bfseries\color{bidaccent}4 · Fuentes de busqueda}\\[5pt]
{\color{bidgray}\normalsize 9 fuentes · 5 libres · 2 con key · 2 pendientes}
\end{center}
\end{frame}
}

\begin{frame}{9 fuentes de busqueda academica}{}
\renewcommand{\arraystretch}{1.2}
\setlength\tabcolsep{4pt}
\footnotesize
\begin{tabular}{@{}ll>{\raggedright\arraybackslash}p{58mm}l@{}}
\toprule
\textbf{Fuente} & \textbf{Auth} & \textbf{Mejor para} & \textbf{Estado} \\
\midrule
Semantic Scholar & libre   & CS, interdisciplinar, ArXiv+journals          & \statuspill{activa} \\
NBER             & libre   & WPs economia: Brynjolfsson, Autor, Acemoglu   & \statuspill{activa} \\
OpenAlex         & libre   & 240M works, citation counts verificados       & \statuspill{activa} \\
ArXiv            & libre   & Preprints CS/AI 2024--2026                    & \statuspill{activa} \\
World Bank       & libre   & WPs desarrollo, labor, LATAM                  & \statuspill{activa} \\
CORE             & Bearer  & Texto completo open access                    & \statuspill{activa} \\
SSRN             & Elsevier& Preprints econ/finanzas                       & \statuspill[colamber]{pendiente} \\
PubMed           & NCBI    & Biomed (solo si pocket = salud)               & \statuspill[bidmid]{auto-skip} \\
IMF              & propia  & WPs macroeconomia                             & \statuspill[bidmid]{sin key} \\
\bottomrule
\end{tabular}

\vspace{3mm}
\begin{columns}[T]
\column{0.48\textwidth}
\begin{card}
\scriptsize\color{bidgray}
\textbf{\color{bidaccent}Crossref:} \texttt{enrich\_metadata} agrega journal, ISSN, citas formales y ORCID por DOI. Sin key, sin costo.
\end{card}
\column{0.48\textwidth}
\begin{card}
\scriptsize\color{bidgray}
\textbf{\color{bidaccent}Resultados guardados 24h:} segunda corrida del mismo pocket = 0 llamadas a APIs de busqueda.
\end{card}
\end{columns}
\end{frame}

% ══════════════════════════════════════════════════════════════
% 5. SCORING
% ══════════════════════════════════════════════════════════════
{
\setbeamercolor{background canvas}{bg=bidmid}
\begin{frame}[plain]\vspace{20mm}
\begin{center}
{\Large\bfseries\color{bidaccent}5 · Sistema de scoring v2}\\[5pt]
{\color{bidgray}\normalsize Inspirado en GRADE + EPPI-Centre Weight of Evidence · escala 0--14}
\end{center}
\end{frame}
}

\begin{frame}{Sistema de scoring v2}{5 dimensiones independientes · escala 0--14}
\renewcommand{\arraystretch}{1.4}
\setlength\tabcolsep{4pt}
\small
\begin{tabular}{@{}p{34mm}cp{84mm}@{}}
\toprule
\textbf{Dimension} & \textbf{Pts} & \textbf{Descripcion} \\
\midrule
D1 Identificacion causal  & 0--4 &
  4=RCT preregistrado \enspace 3=DiD/IV/RD con tests \enspace 2=Panel FE \enspace 1=OLS \enspace 0=descriptivo \\[2pt]
D2 Seniales de calidad    & 0--3 &
  +1 venue top-5 econ o top CS \enspace +1 citas $>$100 o h-index $>$30 \enspace +1 replicado \\[2pt]
D3 Validez externa LATAM  & 0--3 &
  3=evidencia directa LATAM \enspace 2=MiPyMEs USA/Eur. \enspace 1=firmas grandes \enspace 0=no transferible \\[2pt]
\textbf{D4 Relevancia}    & 0--3 &
  \textbf{Gate absoluto: D4=0 implica RECHAZADO sin importar D1--D3} \\[2pt]
D5 Recencia ajustada      & $-1$/0/$+1$ &
  $+1$=2023--26 GenAI \enspace 0=2020--22 o fundacional ($>$50 citas) \enspace $-1$=pre-2020 y $<$50 citas \\
\bottomrule
\end{tabular}

\vspace{4mm}
\begin{columns}[T]
\column{0.32\textwidth}
\begin{tcolorbox}[colback=colgreen!20!cardbg,colframe=colgreen!60,arc=3pt,
  boxrule=0.5pt,left=4pt,right=4pt,top=3pt,bottom=3pt]
\scriptsize\centering
\textbf{\color{colgreen!90}ACEPTADO}\\[3pt]
\color{bidgray}score $\geq$ 9 \\ o bien score $\geq$7 con D4=3 y D1$\geq$2
\end{tcolorbox}
\column{0.32\textwidth}
\begin{tcolorbox}[colback=colamber!20!cardbg,colframe=colamber!60,arc=3pt,
  boxrule=0.5pt,left=4pt,right=4pt,top=3pt,bottom=3pt]
\scriptsize\centering
\textbf{\color{colamber!90}REVISAR}\\[3pt]
\color{bidgray}score 5--8\\revision manual
\end{tcolorbox}
\column{0.32\textwidth}
\begin{tcolorbox}[colback=colred!20!cardbg,colframe=colred!60,arc=3pt,
  boxrule=0.5pt,left=4pt,right=4pt,top=3pt,bottom=3pt]
\scriptsize\centering
\textbf{\color{colred!80}RECHAZADO}\\[3pt]
\color{bidgray}score $<$ 5 \\ o D4 = 0
\end{tcolorbox}
\end{columns}
\end{frame}

\begin{frame}{Sistema de scoring v2}{Frameworks de referencia: de donde viene cada dimension}
\begin{columns}[T]
\column{0.32\textwidth}
\begin{card}
\small\textbf{\color{bidaccent}GRADE}\\
{\tiny\color{bidaccent!60}¿Lo usamos? Si, adaptado}\\[4pt]
\scriptsize\color{bidgray}
Estandar en medicina y economia del desarrollo. Separa \emph{calidad interna} del estudio de su \emph{aplicabilidad} a politica.\\[4pt]
\textbf{Como lo adaptamos:} la jerarquia RCT $>$ DiD/IV $>$ Panel $>$ OLS $>$ descriptivo viene directamente de GRADE. D1 es esencialmente su escala de evidencia causal.\\[4pt]
\textbf{Que no tomamos:} GRADE tiene 8 criterios de downgrade (riesgo de sesgo, imprecision, etc.) que no implementamos por costo computacional.
\end{card}

\column{0.32\textwidth}
\begin{card}
\small\textbf{\color{bidaccent}EPPI-Centre WoE}\\
{\tiny\color{bidaccent!60}¿Lo usamos? Si, como estructura}\\[4pt]
\scriptsize\color{bidgray}
Weight of Evidence para revisiones sistematicas en ciencias sociales. Distingue \emph{solidez del estudio} de \emph{pertinencia de la pregunta}.\\[4pt]
\textbf{Como lo adaptamos:} esa distincion es exactamente D1--D2 (solidez) vs.\ D4 (pertinencia). Un buen estudio en EEUU con D1=4 puede tener D4=1 si la pregunta no aplica al pocket.\\[4pt]
\textbf{Que no tomamos:} WoE usa sintesis narrativa manual; nosotros la automatizamos con Haiku.
\end{card}

\column{0.32\textwidth}
\begin{card}
\small\textbf{\color{bidaccent}3ie Evidence Gap Maps}\\
{\tiny\color{bidaccent!60}¿Lo usamos? Para los gaps}\\[4pt]
\scriptsize\color{bidgray}
Metodologia del BID para mapear evidencia disponible vs.\ ausente. Matriz: diseno de identificacion × relevancia tematica. Celdas vacias = gaps.\\[4pt]
\textbf{Como lo adaptamos:} el reporte de gaps que genera Sonnet sigue esta logica: cruza D1 (metodologia) × D3 (contexto LATAM) y reporta donde no hay papers.\\[4pt]
\textbf{Que no tomamos:} 3ie requiere busquedas exhaustivas de literatura gris; nosotros limitamos a fuentes academicas top-tier.
\end{card}
\end{columns}

\vspace{3mm}
\begin{card}
\scriptsize\color{bidgray}
\textbf{\color{bidaccent}En sintesis:} no implementamos ninguno al pie de la letra. GRADE y EPPI informan la \emph{estructura} de D1--D4; 3ie informa el \emph{reporte de gaps}. Las tres son reconocidas por el BID, lo que facilita la validacion con contrapartes institucionales.
\end{card}
\end{frame}

% ══════════════════════════════════════════════════════════════
% 6. RESULTADOS
% ══════════════════════════════════════════════════════════════
{
\setbeamercolor{background canvas}{bg=bidmid}
\begin{frame}[plain]\vspace{20mm}
\begin{center}
{\Large\bfseries\color{bidaccent}6 · Resultados}\\[5pt]
{\color{bidgray}\normalsize Primera iteracion --- Pocket de Impacto Causal}
\end{center}
\end{frame}
}

\begin{frame}{Resultados: Pocket de Impacto Causal}{72 papers procesados · primera corrida completa}

% KPI row
\begin{columns}[T]
\column{0.24\textwidth}
\begin{tcolorbox}[colback=cardbg,colframe=bidaccent!50,arc=4pt,boxrule=0.5pt,
  left=4pt,right=4pt,top=4pt,bottom=4pt,halign=center]
\scriptsize\color{bidgray}Total procesados\\[3pt]
{\Large\bfseries\color{white}72}
\end{tcolorbox}
\column{0.24\textwidth}
\begin{tcolorbox}[colback=colgreen!15!cardbg,colframe=colgreen!60,arc=4pt,boxrule=0.5pt,
  left=4pt,right=4pt,top=4pt,bottom=4pt,halign=center]
\scriptsize\color{bidgray}Aceptados\\[3pt]
{\Large\bfseries\color{colgreen!90}26}\\
\tiny\color{bidgray}36\% del total
\end{tcolorbox}
\column{0.24\textwidth}
\begin{tcolorbox}[colback=colamber!15!cardbg,colframe=colamber!60,arc=4pt,boxrule=0.5pt,
  left=4pt,right=4pt,top=4pt,bottom=4pt,halign=center]
\scriptsize\color{bidgray}Por revisar\\[3pt]
{\Large\bfseries\color{colamber}40}\\
\tiny\color{bidgray}56\% del total
\end{tcolorbox}
\column{0.24\textwidth}
\begin{tcolorbox}[colback=colred!10!cardbg,colframe=colred!50,arc=4pt,boxrule=0.5pt,
  left=4pt,right=4pt,top=4pt,bottom=4pt,halign=center]
\scriptsize\color{bidgray}Rechazados\\[3pt]
{\Large\bfseries\color{colred!80}6}\\
\tiny\color{bidgray}8\% del total
\end{tcolorbox}
\end{columns}

\vspace{4mm}
\begin{columns}[T]
\column{0.48\textwidth}
\begin{card}
\small\textbf{\color{bidaccent}Metricas clave}\\[3pt]
\scriptsize\color{bidgray}
\renewcommand{\arraystretch}{1.2}
\setlength\tabcolsep{4pt}
\begin{tabular}{@{}lr@{}}
Coverage score           & \textbf{0.66} $>$ 0.6 \\
Score promedio aceptados & $\approx$10.2 / 14 \\
D1 $\geq$ 3 en aceptados & 88\% \\
D4 = 3 en aceptados      & 100\% \\
Corrida total            & $\approx$22 min \\
Costo API                & $\approx$\$0.12 USD \\
\end{tabular}
\end{card}

\column{0.48\textwidth}
\begin{card}
\small\textbf{\color{bidaccent}Composicion de los 40 REVISAR}\\[3pt]
\scriptsize\color{bidgray}
\begin{itemize}\setlength\itemsep{3pt}
  \item Mayoria: D1$\geq$2 y D4=3 pero D3 bajo (EEUU/Europa)
  \item Son candidatos directos a aceptar tras revision manual
  \item Potencial: corpus puede subir a $\approx$55--60 aceptados
\end{itemize}
\vspace{3pt}
\textbf{\color{bidaccent}Fuentes con mayor yield:} NBER (Brynjolfsson, Humlum), Semantic Scholar (Dell'Acqua, Noy \& Zhang), ArXiv preprints 2024--25
\end{card}
\end{columns}
\end{frame}

\begin{frame}{Gaps identificados por el agente}{Pocket de Impacto Causal · output del reporte de Sonnet}
\begin{columns}[T]
\column{0.48\textwidth}
\begin{card}
\small\textbf{\color{bidaccent}Gaps de cobertura detectados}\\[5pt]
\scriptsize\color{bidgray}
\begin{itemize}\setlength\itemsep{5pt}
  \item \textbf{LATAM / Sur Global:} $<$5\% del corpus tiene evidencia directa de LATAM. Mayoria son papers de EEUU/Europa.
  \item \textbf{MiPyMEs:} los RCTs disponibles son con trabajadores de grandes firmas (call centers, consultoras, tech). Evidencia en pequenias firmas es escasa.
  \item \textbf{2024--2026:} cobertura mejora con ArXiv pero muchos papers de vanguardia son preprints sin revision formal.
  \item \textbf{Sectores:} concentracion en knowledge workers y CS. Manufactura y servicios basicos poco representados.
\end{itemize}
\end{card}

\column{0.48\textwidth}
\begin{card}
\small\textbf{\color{bidaccent}Implicaciones para el RCT Guatiguara}\\[5pt]
\scriptsize\color{bidgray}
\begin{itemize}\setlength\itemsep{5pt}
  \item El gap LATAM/MiPyME \textbf{es la razon de ser del experimento}: la literatura global no responde las preguntas locales.
  \item Heterogeneidad \textbf{novato vs.\ experto} bien documentada en EEUU --- hipotesis directamente transferible al diseno.
  \item Outcomes mas replicables: productividad por hora y calidad de output. Son los mas medidos globalmente.
  \item Falta evidencia sobre \textbf{efectos de mediano plazo} ($>$6 meses): oportunidad de contribucion original del RCT.
\end{itemize}
\end{card}
\end{columns}
\end{frame}

% ══════════════════════════════════════════════════════════════
% 7. PRÓXIMOS PASOS
% ══════════════════════════════════════════════════════════════
{
\setbeamercolor{background canvas}{bg=bidmid}
\begin{frame}[plain]\vspace{20mm}
\begin{center}
{\Large\bfseries\color{bidaccent}7 · Proximos pasos}\\[5pt]
{\color{bidgray}\normalsize Agenda para la siguiente iteracion}
\end{center}
\end{frame}
}

\begin{frame}{Proximos pasos}{Dos rutas en paralelo}
\begin{columns}[T]
\column{0.48\textwidth}
\begin{card}
\small\textbf{\color{bidaccent}Pipeline (automatico)}\\[5pt]
\scriptsize\color{bidgray}
\begin{enumerate}\setlength\itemsep{5pt}
  \item Correr \texttt{--mode search} en los 6 pockets restantes \\ estimado \$3--4 USD
  \item Correr \texttt{build\_network.py} para grafo de citaciones \\ $\approx$\$0.05 USD
  \item Actualizar \texttt{main\_dashboard.html} con corpus completo
  \item Resolver SSRN API key con Elsevier
\end{enumerate}
\end{card}

\column{0.48\textwidth}
\begin{card}
\small\textbf{\color{bidaccent}Revision manual (investigador)}\\[5pt]
\scriptsize\color{bidgray}
\begin{enumerate}\setlength\itemsep{5pt}
  \item Revisar los \textbf{40 papers en REVISAR} del Pocket de Impacto Causal via \texttt{review\_interface.html}
  \item Validar los 26 aceptados contra el juicio del equipo
  \item Ajustar umbrales de scoring si hay discrepancias sistematicas
  \item Agregar manualmente anchor papers no encontrados por las APIs
\end{enumerate}
\end{card}
\end{columns}

\vspace{5mm}
\begin{center}
\begin{tikzpicture}
\node[rectangle,rounded corners=3pt,fill=bidlight,text=white,
  font=\small\bfseries,minimum width=110mm,minimum height=9mm,inner sep=4pt]
  {Meta DoD: $>$50 aceptados por pocket · Impacto Causal: 26 confirmados + 40 por revisar};
\end{tikzpicture}
\end{center}
\end{frame}

% ── Anexo ────────────────────────────────────────────────────
{
\setbeamercolor{background canvas}{bg=bidmid}
\begin{frame}[plain]\vspace{20mm}
\begin{center}
{\Large\bfseries\color{bidaccent}Anexo}\\[5pt]
{\color{bidgray}\normalsize Ejemplo comparativo de frameworks de clasificacion}
\end{center}
\end{frame}
}

\begin{frame}{Anexo · Frameworks originales}{Estructura real de cada uno}
\begin{columns}[T]

\column{0.32\textwidth}
\begin{card}
\scriptsize\textbf{\color{bidaccent}GRADE}\\
{\tiny\color{bidaccent!60}Grading of Recommendations Assessment}\\[3pt]
\tiny\color{bidgray}
Tabla de certeza por outcome; columnas evaluan downgrade/upgrade.
\vspace{2pt}
\setlength\tabcolsep{2pt}
\renewcommand{\arraystretch}{1.15}
\begin{tabular}{@{}p{11mm}ccc@{}}
\toprule
\tiny\textbf{Outcome} & \tiny\textbf{Diseno} & \tiny\textbf{Sesgo} & \tiny\textbf{Cert.} \\
\midrule
Velocidad & RCT & Bajo & \statuspill{Alta} \\
Calidad & RCT & Bajo & \statuspill{Alta} \\
LATAM & — & — & \statuspill[colred]{Baja} \\
\bottomrule
\end{tabular}
\end{card}

\column{0.32\textwidth}
\begin{card}
\scriptsize\textbf{\color{bidaccent}EPPI-Centre WoE}\\
{\tiny\color{bidaccent!60}Weight of Evidence · 4 componentes}\\[3pt]
\tiny\color{bidgray}
Tres pesos independientes combinados en veredicto WoE D.
\vspace{2pt}
\setlength\tabcolsep{2pt}
\renewcommand{\arraystretch}{1.15}
\begin{tabular}{@{}p{9mm}p{16mm}c@{}}
\toprule
\tiny\textbf{Comp.} & \tiny\textbf{Pregunta} & \tiny\textbf{Peso} \\
\midrule
WoE A & \tiny Calidad metodologica & \statuspill{Alto} \\
WoE B & \tiny Relevancia diseno & \statuspill{Alto} \\
WoE C & \tiny Relevancia foco & \statuspill[colamber]{Medio} \\
\midrule
\textbf{WoE D} & \textbf{Veredicto} & \statuspill{Alto} \\
\bottomrule
\end{tabular}
\end{card}

\column{0.32\textwidth}
\begin{card}
\scriptsize\textbf{\color{bidaccent}3ie Evidence Gap Map}\\
{\tiny\color{bidaccent!60}Matriz intervencion × outcome}\\[3pt]
\tiny\color{bidgray}
Cada celda = evidencia disponible. Celda vacia = gap.
\vspace{2pt}
\setlength\tabcolsep{2pt}
\renewcommand{\arraystretch}{1.15}
\begin{tabular}{@{}p{11mm}ccc@{}}
\toprule
 & \tiny Productiv. & \tiny Empleo & \tiny Bienestar \\
\midrule
IA · EEUU/Eur & \statuspill{Alto} & \statuspill[colamber]{Med} & \statuspill[colamber]{Bajo} \\
IA · LATAM    & \statuspill[colamber]{Bajo} & \statuspill[colred]{Gap} & \statuspill[colred]{Gap} \\
IA · MiPyMEs  & \statuspill[colred]{Gap} & \statuspill[colred]{Gap} & \statuspill[colred]{Gap} \\
\bottomrule
\end{tabular}
\end{card}

\end{columns}

\vspace{3mm}
\begin{card}
\tiny\color{bidgray}
\renewcommand{\arraystretch}{1.2}
\setlength\tabcolsep{4pt}
\begin{tabular}{@{}p{18mm}p{55mm}p{55mm}@{}}
 & \textbf{\color{bidaccent}Que tomamos} & \textbf{\color{bidaccent}Que dejamos} \\
\midrule
\textbf{GRADE}  & Jerarquia causal RCT$>$DiD$>$Panel$>$OLS (D1) & 8 criterios de downgrade por riesgo de sesgo interno \\
\textbf{EPPI}   & Separacion calidad del estudio (D1--D2) vs.\ pertinencia (D4) & Sintesis narrativa manual por panel de expertos \\
\textbf{3ie}    & Matriz gaps intervencion × contexto → reporte de Sonnet & Busqueda exhaustiva de literatura gris y no publicada \\
\end{tabular}
\end{card}
\end{frame}

% ── Cierre ────────────────────────────────────────────────────
{
\setbeamercolor{background canvas}{bg=bidblue}
\begin{frame}[plain]
\vspace{14mm}\hspace{7mm}%
\begin{minipage}{0.76\textwidth}
{\Large\bfseries\color{white}Gracias}
\vspace{6mm}

\textcolor{bidgray}{\small
El codigo, los datos y los dashboards interactivos estan\\disponibles en el repositorio del proyecto.}

\vspace{8mm}
\begin{card}
\scriptsize\color{bidgray}
\texttt{\$ python scripts/agent.py --pocket evaluacion\_experimental --mode full}\\[3pt]
\textcolor{bidaccent!70}{Pipeline completo: busqueda · enriquecimiento · deduplicacion · review · sintesis · guardado}\\
\textcolor{bidaccent!70}{Costo: $\approx$\$0.12 USD · Tiempo: $\approx$20 min}
\end{card}

\vspace{6mm}
\textcolor{bidgray!50}{\footnotesize
  Maria Alejandra Montoya · m.montoyac@uniandes.edu.co\\
  Fedesarrollo \& BID · Mayo 2026}
\end{minipage}
\end{frame}
}

\end{document}